\section{Discussion}
\label{sec:discussion}

\subsection{Hypothesis Assessment}

\tabref{tab:hyp} maps our predictions to observations.

\begin{table}[t]
\centering
\caption{Hypothesis assessment. ``partial'' = direction observed in some
conditions but not others; ``flipped'' = predicted direction not observed.}
\label{tab:hyp}
\begin{tabular}{@{}llll@{}}
\toprule
\textbf{Sub-hypothesis} & \textbf{Predicted} & \textbf{Observed} & \textbf{Verdict} \\
\midrule
H1: \layernorm compresses concept distance      & yes  & \pythia: $+7\%$; \gptsmall: $\sim 0$ & partial \\
H2: ablating \layernorm increases \dat          & yes  & only at high-$T$ sampling          & partial \\
H3: effect localised to early layers            & yes  & uniform across targets at high-$T$ & flipped \\
H4: effect persists across decoders             & yes  & decoder-dependent                  & no \\
H5: surprisal does not solely explain \dat shift & --   & under-tested                       & open \\
\bottomrule
\end{tabular}
\end{table}

\subsection{Why Does the Latent Expansion Not Propagate?}

Three reasons stand out for why the $+7\%$ pair-distance expansion in
\pythia does not surface as a $+7\%$ \dat shift.

\para{The LM head is tuned to the stock geometry.}
The unembedding matrix is a learned linear map onto the vocabulary; its
rows are aligned with the \emph{stock} hidden-state distribution. When
\nomean re-spreads the final hidden states, the LM head's row alignment is
unchanged, so per-vocabulary logits and the resulting softmax change only
modestly. The dominant geometric constraint is the LM head, which is
preserved by all our interventions.

\para{Decoding recompresses the output distribution.}
The \dat prompt is heavily anchored (``the next noun in a list of
unrelated single-word nouns''). The conditional distribution over candidate
tokens is narrow --- only nouns survive top-$p{=}0.9$ at $T{=}0.7$. A few
percent of latent shift does not change the argmax token, and therefore
does not change \dat. At $T{=}1.2$, top-$p{=}0.95$, the candidate set is
broader, and small probability shifts can flip the sampled token. This is
why the high-$T$ panel is the only one where the geometric expansion
surfaces in the behavioural metric.

\para{The final block is preserved.}
Following \citet{kanavalau2026taper}, we leave the final \layernorm intact
to keep the LM head's input distribution calibrated. This final cleanup
absorbs much of the intermediate variation introduced by upstream
ablations.

\subsection{Alignment with Prior Work}

\para{Geometry.}
\citet{brody2023expressivity} and \citet{gupta2024geometric} predict that
\nomean is near-redundant on already-zero-mean hidden states. We
reproduce this: \nomean changes perplexity by $3\!-\!12\%$ and the angle
to $\bm{1}$ stays $\sim\!88\!-\!92^\circ$ regardless of intervention.

\para{Subspace coupling.}
\citet{menary2024transformer} predict that O($10\%$) L2-norm perturbations
collapse $\sim\!1\%$ of attention heads. We observe \emph{catastrophic}
breakage at $\alpha\le 0.5$ \weakened (PPL $> 1{,}000$) --- consistent with
crossing the threshold for many heads simultaneously.

\para{Layer specificity.}
\citet{singhal2025impact} find early-layer \layernorm most influential for
training-time memorisation. Our \dat data does \emph{not} show
early-layer specificity at inference. The discrepancy may reflect that
their effects are locked in early during training and are not easily
perturbed by inference-time hooks.

\para{Internal vs.\ output randomness.}
\citet{nagarajan2025roll} argue that internal randomness beats output
randomness for diversity. Our results do not contradict this, but suggest
\layernorm is not the right place to inject internal randomness:
input-level seed conditioning, as they propose, likely engages different
mechanisms.

\subsection{Surprises}

Three results surprised us.
First, \tgtlate \nomean is the \emph{biggest} (positive but
non-significant) effect in \gptsmall \topp --- the \emph{opposite} of the
\citet{singhal2025impact} early-layer memorisation story. Second,
broken interventions inflate \dat \emph{artifactually} via gibberish words
with random embeddings; we are not aware of prior \dat-LLM work that gates
on validity rate as we do, and we recommend it as a methodological
default. Third, \pythia and \gptsmall disagree on the geometric
\emph{direction} of \layernorm's effect at the final layer: \pythia
expands under \nomean, \gptsmall contracts. This hints at architectural
or training-data-specific behaviour and motivates multi-model evaluation
of any \layernorm-creativity intervention.

\subsection{Limitations}

\para{Inference-time only.}
We do not retrain. \textsc{TaperNorm} \citep{kanavalau2026taper}, \textsc{nGPT}
\citep{loshchilov2025ngpt}, and \citet{singhal2025impact} all show
training-time effects of \layernorm that dominate. A fully causal test
would compare models trained with vs.\ without \layernorm on the same
data. This is beyond a single-session compute budget but is the natural
follow-up.

\para{Small models.}
We study models in the $124\!-\!410$M range. The concept-incongruence
phenomenon is most observable in much larger instruction-tuned models;
generalisation to \textsc{Llama}-class is open.

\para{Final \layernorm preserved.}
We intentionally keep the last \layernorm so the LM head's input
distribution stays calibrated. Ablating that final \layernorm produces
uninterpretable outputs and is a different experiment.

\para{One creativity metric.}
\dat is one creativity probe; algorithmic-creativity tasks
\citep{nagarajan2025roll}, the Alternative Uses Test, and creative writing
were not tested. The \dat-specific null is not a global null on
creativity-vs-\layernorm.

\para{Iterative one-noun protocol.}
\citet{chen2023probing} use one-shot list generation on
instruction-tuned LLMs; base \gptsmall cannot. Our iterative protocol asks
for one noun given the partial list each step, which conditions on a
growing unrelated context. Absolute \dat numbers should be interpreted
under this protocol; relative comparisons across interventions are
within-protocol and valid.

\subsection{Broader Implications}

The widespread intuition --- that softening \layernorm at inference might
unlock creative diversity --- is not borne out by our experiments.
\layernorm's geometric grip on hidden states is real and reproducible, but
it is not the dominant gate between latent variation and output diversity.
The LM-head and softmax recompress most of the freed variation, and only
when the decoder is also relaxed (\eg high temperature) does any of it
surface in \dat. For practitioners interested in creative-divergence
gains, the cleaner levers are decoder design (min-$p$ sampling
\citep{nguyen2025minp}; seed-conditioning \citep{nagarajan2025roll}) and
training-time interventions (\textsc{TaperNorm}, \textsc{nGPT}). For
mechanistic interpretability, our results recommend that any future
\layernorm-creativity claim be co-validated against (i) a fluency check,
(ii) a validity-rate gate, and (iii) more than one decoder regime.
